% \documentclass[polish, 8pt]{article}
\documentclass[polish, 10pt]{extarticle}
% \usepackage[top=2cm, bottom=2cm, left=3cm, right=3cm]{geometry}
\usepackage[a4paper,top=1.2cm,bottom=1.2cm,left=1.2cm,right=1.2cm]{geometry}
\usepackage{graphicx} % Required for inserting images

\usepackage[polish]{babel}
\usepackage{polski}
\usepackage{placeins}
\usepackage{dirtytalk}
% Set page size and margins
% Replace `letterpaper' with `a4paper' for UK/EU standard size
\usepackage{csquotes}
% Useful packages
\usepackage{float}
\usepackage{amsmath}
\usepackage{graphicx}
\usepackage{hyperref}
\usepackage{graphicx}
\usepackage{listings}
\usepackage{xcolor}
\usepackage{amssymb}
\usepackage{enumitem}
\usepackage{float}
\definecolor{codegreen}{rgb}{0,0.6,0}
\definecolor{codegray}{rgb}{0.5,0.5,0.5}
\definecolor{codepurple}{rgb}{0.58,0,0.82}
\definecolor{backcolour}{rgb}{0.95,0.95,0.92}

\lstdefinestyle{mystyle}{
    backgroundcolor=\color{backcolour},   
    commentstyle=\color{codegreen},
    keywordstyle=\color{magenta},
    numberstyle=\tiny\color{codegray},
    stringstyle=\color{codepurple},
    basicstyle=\ttfamily\footnotesize,
    breakatwhitespace=false,         
    breaklines=true,                 
    captionpos=b,                    
    keepspaces=true,                 
    numbers=left,                    
    numbersep=5pt,                  
    showspaces=false,                
    showstringspaces=false,
    showtabs=false,                  
    tabsize=2
}
\lstset{style=mystyle}
\hypersetup{
    colorlinks=true,
    linktoc=all     %set to all if you want both sections and subsections linked
}
\usepackage{mdframed}
\usepackage{titling}
\setlength{\droptitle}{-2cm} 

\begin{document}

\title{Kolokwium II -- Zestaw B}
\author{}
\date{}

\maketitle
\vspace{-2cm}

\section*{Zadanie 1 (30 pkt)}
Niech dwuwymiarowa zmienna losowa $(X, Y)$ ma rozkład ciągły o gęstości:
\[
f(x,y) = 
\begin{cases} 
cx & \text{dla } 0 \le x \le 1, \,\, x^2 \le y \le 1, \\
0 & \text{w przeciwnym przypadku,}
\end{cases}
\]
gdzie $c \in \mathbb{R_+}$ jest pewną stałą.

\begin{enumerate}
    \item[a)] \textbf{Wyznacz wartość stałej $c$} \hfill (5 pkt).
    \item[b)] \textbf{Wyznacz gęstości brzegowe} \hfill (5 pkt).
    \item[c)] \textbf{Wyznacz gęstości warunkowe} \hfill (5 pkt).
    \item[d)] \textbf{Wyznacz $\mathbb P\!\left(Y>\tfrac12\mid X=\tfrac12\right)$} \hfill (5 pkt).
    \item[e)] \textbf{Oblicz kowariancję $\text{Cov}(X, Y)$} \hfill (5 pkt).
    \item[f)] \textbf{Opisz sposób wyznaczenia współczynnika korelacji $\rho_{X,Y}$},\\rozpisując wzory i całki niezbędne do jego obliczenia \hfill (5 pkt).
\end{enumerate}



\section*{Zadanie 2 (15 pkt)}
Dla zmiennej losowej $X\sim \mathrm{Pois}(\lambda)$:
\begin{enumerate}
  \item[a)] Wyznacz funkcję tworzącą momenty (MGF) \hfill (10 pkt).
  \item[b)] Korzystając z MGF oblicz wariancję $\mathrm{Var}(X)$ \hfill (5 pkt).
\end{enumerate}

\section*{Zadanie 3 (30 pkt)}
System zasilania awaryjnego w serwerowni składa się z $n$ niezależnych modułów bateryjnych. Moduły te włączają się sekwencyjnie – gdy pierwszy się wyczerpie, natychmiast startuje drugi, itd. Czas pracy każdego $i$-tego modułu (w godzinach) jest zmienną losową $X_i$ o rozkładzie wykładniczym z parametrem $\lambda = 2$: $X_i \sim \text{Exp}(2)$. Niech $T$ oznacza całkowity czas pracy systemu do momentu wyczerpania wszystkich baterii:
\[
T = X_1 + X_2 + \dots + X_n
\]

\begin{enumerate}
    \item[a)] Oblicz wartość oczekiwaną $\mathbb{E}[T]$ i wariancję $\text{Var}(T)$ \hfill (5 pkt).
    \item[b)] Korzystając z własności sumy niezależnych zmiennych losowych, podaj dokładny rozkład zmiennej $T$ \hfill (5 pkt).
    \item[c)] Załóżmy, że system składa się z $n=100$ modułów. Oszacuj prawdopodobieństwo, że system będzie działał dłużej niż 55 godzin. Powołaj się na odpowiednie twierdzenie \hfill (10 pkt).
    \item[d)] Traktując parametr $\lambda$ jako nieznany, a czasy pracy modułów $x_1,\dots,x_n$ jako zaobserwowane realizacje niezależnych zmiennych $X_1,\dots,X_n \sim \mathrm{Exp}(\lambda)$, wyznacz estymator największej wiarygodności $\hat\lambda_{\text{MLE}}$. \hfill (10 pkt)
\end{enumerate}

\section*{Zadanie 4 (25 pkt)}
W pewnym zakładzie produkcyjnym kontroluje się czas (w minutach) potrzebny na wykonanie standardowej operacji. Przyjmuje się, że czasy wykonywania tej operacji w populacji mają rozkład normalny. Wylosowano próbę \textbf{$n = 14$} pomiarów, dla których uzyskano $\bar{x} = 12{,}8$ min przy $s = 3{,}2$ min.

\begin{enumerate}
    \item \textbf{Opisz sposób konstrukcji $95\%$ przedziału ufności dla $\mu$}, tj.:
    \begin{itemize}
        \item jaka statystyka jest wykorzystywana do budowy przedziału i jaki jest jej rozkład, \hfill (5 pkt),
        \item dlaczego wybieramy akurat tę statystykę w tej sytuacji, \hfill (5 pkt),
        \item opisz sposób konstrukcji (jeśli wprowadzasz zmienne, wyjaśnij  co oznaczają) \hfill (10 pkt).
    \end{itemize}

    \item Ktoś interpretuje wynik następująco:
    \begin{quote}
    ``Z prawdopodobieństwem $95\%$ prawdziwa wartość średniej $\mu$ leży w wyznaczonym przedziale.''
    \end{quote}
    Czy ta interpretacja jest \textbf{prawdziwa}? Uzasadnij. \hfill (5 pkt).
\end{enumerate}



\end{document}

